% TeX's null delimiter: `\left( ... \right.` is a fence whose right delimiter is
% empty. Digestion emits no token for the `.` (Perl does the same — the XMath is
% byte-identical), so a fence split across an alignment break leaves each cell
% with an unmatched delimiter. The grammar's fence rules for general
% expressions are balanced-only, so this used to yield ZERO parses and drop the
% whole formula — fractions and subscripts included — to ltx_math_unparsed.
% Witness: arXiv 2606.13010 appendix (CR Hamiltonian).
\documentclass{article}
\usepackage{amsmath}
\begin{document}

% The split fence itself: both cells must parse, and neither may gain a
% delimiter the author did not write.
\begin{align}
  H & = \frac{\hbar}{2} \left( a + b \right. \nonumber \\
    & \quad \left. + c + d \right).
\end{align}

% Balanced control — must be completely unaffected by the repair.
\begin{equation}
  H = \frac{\hbar}{2} \left( a + b + c + d \right)
\end{equation}

% REGRESSION GUARD. A ket's closing `⟩` pairs with a VERTBAR, not with an OPEN,
% so a naive "unbalanced ⇒ synthesize a partner" pass reads it as a dangling
% close and prepends a bogus `(`. These parsed before the repair existed and
% must keep parsing; an earlier unconditional version broke exactly these.
\(|f\rangle\) and \(\langle g|\) and \(\langle f | g \rangle\)

\end{document}
